\documentclass{book}
\usepackage[utf8]{inputenc}
\usepackage[T1]{fontenc}
\usepackage{geometry}
\usepackage{graphicx}
\usepackage{float}
\usepackage{hyperref}
\usepackage[english,indonesian]{babel}

\title{Optimasi Pemrosesan Template LaTeX Menggunakan Pendekatan Berbasis AI}
\author{Academic Support Agent}
\date{\today}

\begin{document}
\maketitle

\begin{abstract}

\end{abstract}

\chapter{PENDAHULUAN}


\chapter{TINJAUAN PUSTAKA}


% Original content (commented out):
% \title{Διδακτορική Διατριβή}
\author{Λεοντιάδης Κωνσταντίνος}
\date{Θεσσαλονίκη 2023}
%====================================
%====================================

% Η σύνταξη του κειμένου γίνεται με τον compiler XeLaTeX.

% Η κλάσση του αρχείου ορίζεται ως book με γραμματοσειρά κυρίως κειμένου μεγέθους 12pt. 
\documentclass[12pt]{book}
\setcounter{tocdepth}{3}
\setcounter{secnumdepth}{3}


% Πακέτα για την ταυτόχρονη χρήση ελληνικών αλλά και αγγλικών χαρακτήρων στο κείμενο.
\usepackage[greek,english, main = greek]{babel} % use this line if greek version is desired
\usepackage{fontspec}
\setmainfont{Libertinus Serif}

% CHECKLATER - Παράδειγμα
% Πακέτο για τον χρωματισμό κειμένου.
\usepackage[dvipsnames]{xcolor}

% Ρυθμίσεις για την διάταξη και τα όρια των σελίδων.
\usepackage[left=3cm, right=2.5cm, top=2.5cm, bottom=2.5cm, includefoot, headheight=13.6pt]{geometry}

% Εσοχή στην πρώτη γραμμή της παραγράφου.
\usepackage{indentfirst}
% και ρύθμιση του μέγεθος της εσοχής της πρώτης γραμμής.
\setlength{\parindent}{25pt}

% CHECKLATER - Παράδειγμα
% Πακέτο για την παράθεση της βιβλιογραφίας σε μορφή IEEE. 
\usepackage[
    backend=biber,
    sorting=none,
    giveninits=true,
    maxbibnames=10,
    bibstyle=ieee,
    citestyle=numeric-comp,
  ]{biblatex}

% Στο αρχείο ref.bib συγκεντρώνονται οι βιβλιογραφικές αναφορές που χρησιμοποιούνται στο κείμενο. 
\addbibresource{ref.bib}

% CHECKLATER
\usepackage{csquotes}

% Εισαγωγή κεφαλίδας.
% Τα στοιχεία που εισάγονται στην κεφαλίδα ρυθμίζονται στο κυρίως κείμενο. 
\usepackage[]{fancyhdr}

% Ώστε η αυτόματη αρίθμηση να σημειώνεται στο κάτω μέρος της σελίδας.
\pagestyle{plain}

% Πακέτα για την εισαγωγή εικόνων
\usepackage{adjustbox}
\usepackage{graphicx}
\usepackage{placeins}
\usepackage{subcaption}
\usepackage{float}

% Πακέτο για να μην ξεκινούν τίτλοι κεφαλαίων στο τέλος σελίδων. 
\usepackage[nobottomtitles*]{titlesec}

% CHECKLATER - Παράδειγμα
% Για την εισαγωγή χημικών τύπων.
\usepackage{chemfig}

% CHECKLATER - Παράδειγμα
% Πακέτο για να είναι "clickable" οι συνδέσεις (references) στο κείμενο
\usepackage[colorlinks=true,citecolor=blue,hidelinks]{hyperref}

% Για τις λεζάντες των στοιχείων του κειμένου
\usepackage[font=small,labelfont=bf]{caption}

% CHECKLATER - Παράδειγμα
% Για την χρήση συγχώνευσης γραμμών στους πίνακες 
\usepackage{multirow}

% CHECKLATER - Παράδειγμα
% Πακέτο για χρήση πιο ευέλικτων πινάκων.
\usepackage{tabularx}
\usepackage{makecell}

% CHECKLATER - Παράδειγμα
% Εντολή για τοποθέτηση σχήματος με λεζάντα 
% Χρήση μέσα στο κείμενο με τη μορφή:
%\addfigure{όνομα_φακέλου}{όνομα_αρχείου}{πλάτος εικόνας}{Λεζάντα εικόνας}
\newcommand{\addfigure}[4]{
\graphicspath{{./#1/}}
 \begin{figure}[H]
 \centering
 \includegraphics[width = #3]{#2}
 \caption{#4}
 \label{#2}
 \end{figure} 
}

% CHECKLATER - Παράδειγμα
% Για την μορφοποίηση όρων μέσα στο κείμενο (εκθέτες, δείκτες κτλ.)
\newcommand{\term}[1]{
$\mathrm{#1}$
}

% CHECKLATER - Παράδειγμα
% Εντολή για προσθήκη πίνακα με λεζάντα 
% Χρήση μέσα στο κείμενο με τη μορφή:
%\addtable{στήλες πχ. "c c c" για τρεις στήλες}{Πρώτη γραμμή πίνακα}{περιεχόμενα πίνακα}{κεφάλαιο πίνακα}{αριθμός πίνακα}{λεζάντα πίνακα}
\newcommand{\addtable}[6]{
 \begin{table}[H]
 \centering
 \caption{#6}
 \begin{tabular}{#1}
 #2
 \hline\hline 
 #3
 \end{tabular}
 \label{tab_#4_#5}
\end{table}
}

% CHECKLATER - Παράδειγμα
% Για σχόλια πολλών σειρών
\usepackage{verbatim}

% CHECKLATER - Παράδειγμα
% Αναφορά σε όνομα κεφαλαίου
\usepackage{nameref}

% CHECKLATER - Παράδειγμα
%Πακέτο για τη χρήση μαθηματικών τύπων.
\usepackage{amsmath}

%=================================================================

\begin{document}

%=========================================
% Παράμετροι για τη μορφοποίηση του κειμένου ώστε να μην ξεφεύγει από τα όρια της σελίδας που έχει τεθεί.
\tolerance 1414
\hbadness 1414
\emergencystretch 1.5em
\hfuzz 0.3pt
\widowpenalty=10000
\vfuzz \hfuzz
\raggedbottom
\linespread{1.25}
%==============================================

% Εντολή για την μη αρίθμηση έως ότου εκτελεστεί η εντολή \frontmatter.
\pagenumbering{gobble} 

% Τα αρχεία τα οποία δεν θα αριθμηθούν βρίσκονται εντός των εντολών \begingroup και \endgroup. 
\begingroup
\include{Cover}
\include{Cover 2}
\include{Chapter -1}
\include{Chapter 0}
\include{acronym catalog}
\include{symbol catalog}
\endgroup 


% Με αυτή την εντολή ξεκινάει η αρίθμηση με μικρούς λατινικούς χαρακτήρες.
\frontmatter

\include{summary}

% Εισαγωγή πίνακα περιεχομένων.
\tableofcontents
% Εισαγωγή κατάλογου σχημάτων.
\listoffigures
% Εισαγωγή κατάλογου πινάκων.
\listoftables


% Αρίθμησή σελίδων με αραβικούς χαρακτήρες (1,2,3 κτλ.)
\mainmatter

% Ο αστερίσκος χρησιμοποιηείτε για να μην αριθμηθεί το κεφάλαιο της Εισαγωγής.
\chapter*{Εισαγωγή}
\label{Introduction}
\input{Introduction}
% Εμφανίζει την εισαγωγή στον πίνακα περιεχομένων (Table of contents -> toc). 
\addcontentsline{toc}{chapter}{Εισαγωγή} 
\cleardoublepage

% Εισαγωγή κεφαλίδας 
\pagestyle{fancy}
\fancyhead{}
\renewcommand{\chaptermark}[1]{\markboth{#1}{#1}}
\fancyhead[RO]{}
\fancyhead[LE]{\chaptername\ \thechapter\ --\ \leftmark}
\setlength\headheight{28pt}

%========================
% Εισαγωγή νέων κεφαλαίων
%========================

\chapter{Οδηγίες}
\label{Chapter 1}
\input{Chapter 1.tex}

\chapter{Δομή}
\label{Chapter 2}
\input{Chapter 2}

\chapter{Παραδείγματα}
\label{Chapter 3}
\input{Chapter 3}


%======================

% Εισαγωγή βιβλιογραφίας

% Αφαίρεση κεφαλίδας για το κεφάλαιο της βιβλιογραφίας. 
\fancyhead[LE]{\chaptername\--\ \leftmark}
\clearpage 

% Εισαγωγή βιβλιγοραφία και ρυθμίσεις για τον τίτλο του κεφαλαίου.
\defbibnote{bibl}{\markboth{Βιβλιογραφία}{Βιβλιογραφία}}
\printbibliography[title={Βιβλιογραφία},prenote=bibl,heading=bibintoc]
\cleardoublepage

% Παράρτηματα
% Επαναφορά κεφαλίδας. 
\fancyhead[LE]{\chaptername\ \thechapter\ --\ \leftmark}
\appendix 

\chapter{Παράρτημα Α}
\label{AppA}
\input{Appendix A}

\chapter{Παράρτημα Κεφαλαίου \ref{Chapter 2}}
\label{AppB}
\input{Appendix B}

\end{document}


\end{document}